% GENERATED FILE. Edit canonical lessons, then run npm run generate:books.
% canonical-source-hash: fnv1a64:4419f3feb96847c4
% canonical-lessons: KA-C193-rattu, KA-C193-akara, KA-C193-vrtta, KA-C193-cauka, KA-C193-trikona, KA-R193-first-pass-words-from-chapters-185-188, KA-R193-second-pass-words-from-chapters-189-193

\chapter{Cardboard, Shape, Circle, Square, Triangle}
\label{ch:ka-cardboard-shape-circle-square-triangle}
\hlchaptermodality{193}

\begin{chapteropening}
\textbf{By the end of this chapter:} \emph{I can say cardboard, a shape, a circle, a square and a triangle.}

Complete the last lesson of chapter 193: I can say cardboard, a shape, a circle, a square and a triangle.
\end{chapteropening}

\section[\texorpdfstring{\kn{ರಟ್ಟು} (raṭṭu)}{ರಟ್ಟು (raṭṭu)}]{\kn{ರಟ್ಟು} (raṭṭu) — cardboard}

\label{lesson:KA-C193-rattu}



\begin{quote}
Before the new one: say the Kannada for a sieve, then the Kannada for a coconut grater.
\end{quote}

\subsection*{You'll want to know: \kn{ರಟ್ಟು}}
\textbf{\kn{ರಟ್ಟು}} — \emph{raṭṭu} — \textquotedblleft{}cardboard\textquotedblright{}.

\begin{grammarlens}[title={What you've built}]
One more word for this chapter.
\end{grammarlens}

\subsection*{Your turn}
\begin{itemize}
\raggedright
  \item \emph{Say it:} \emph{raṭṭu}
  \item \emph{Say it:} \emph{raṭṭu}, once more
  \item \emph{Say it:} say \emph{raṭṭu}
  \item \emph{Recall:} say the Kannada for a sieve, then the Kannada for a coconut grater, then say \emph{raṭṭu} again
\end{itemize}

\begin{tcolorbox}[breakable,colback=teal!4,colframe=teal!35!black,title={Before you move on}]
What does \kn{ರಟ್ಟು} mean? (\textquotedblleft{}Cardboard\textquotedblright{}.) Say it once more.
\end{tcolorbox}
\section[\texorpdfstring{\kn{ಆಕಾರ} (ākāra)}{ಆಕಾರ (ākāra)}]{\kn{ಆಕಾರ} (ākāra) — a shape}

\label{lesson:KA-C193-akara}



\begin{quote}
Before the new one: say the Kannada for a coconut grater, then the Kannada for cardboard.
\end{quote}

\subsection*{You'll want to know: \kn{ಆಕಾರ}}
\textbf{\kn{ಆಕಾರ}} — \emph{ākāra} — \textquotedblleft{}a shape\textquotedblright{}.

\begin{grammarlens}[title={What you've built}]
One more word for this chapter.
\end{grammarlens}

\subsection*{Your turn}
\begin{itemize}
\raggedright
  \item \emph{Say it:} \emph{ākāra}
  \item \emph{Say it:} \emph{ākāra}, once more
  \item \emph{Say it:} say \emph{ākāra}
  \item \emph{Recall:} say the Kannada for a coconut grater, then the Kannada for cardboard, then say \emph{ākāra} again
\end{itemize}

\begin{tcolorbox}[breakable,colback=teal!4,colframe=teal!35!black,title={Before you move on}]
What does \kn{ಆಕಾರ} mean? (\textquotedblleft{}A shape\textquotedblright{}.) Say it once more.
\end{tcolorbox}
\section[\texorpdfstring{\kn{ವೃತ್ತ} (vṛtta)}{ವೃತ್ತ (vṛtta)}]{\kn{ವೃತ್ತ} (vṛtta) — a circle; a traffic circle}

\label{lesson:KA-C193-vrtta}



\begin{quote}
Before the new one: say the Kannada for cardboard, then the Kannada for a shape.
\end{quote}

\subsection*{You'll want to know: \kn{ವೃತ್ತ}}
\textbf{\kn{ವೃತ್ತ}} — \emph{vṛtta} — \textquotedblleft{}a circle; a traffic circle\textquotedblright{}.

\begin{grammarlens}[title={What you've built}]
One more word for this chapter.
\end{grammarlens}

\subsection*{Your turn}
\begin{itemize}
\raggedright
  \item \emph{Say it:} \emph{vṛtta}
  \item \emph{Say it:} \emph{vṛtta}, once more
  \item \emph{Say it:} say \emph{vṛtta}
  \item \emph{Recall:} say the Kannada for cardboard, then the Kannada for a shape, then say \emph{vṛtta} again
\end{itemize}

\begin{tcolorbox}[breakable,colback=teal!4,colframe=teal!35!black,title={Before you move on}]
What does \kn{ವೃತ್ತ} mean? (\textquotedblleft{}A circle; a traffic circle\textquotedblright{}.) Say it once more.
\end{tcolorbox}
\section[\texorpdfstring{\kn{ಚೌಕ} (cauka)}{ಚೌಕ (cauka)}]{\kn{ಚೌಕ} (cauka) — a square}

\label{lesson:KA-C193-cauka}



\begin{quote}
Before the new one: say the Kannada for a shape, then the Kannada for a circle.
\end{quote}

\subsection*{You'll want to know: \kn{ಚೌಕ}}
\textbf{\kn{ಚೌಕ}} — \emph{cauka} — \textquotedblleft{}a square\textquotedblright{}.

\begin{grammarlens}[title={What you've built}]
One more word for this chapter.
\end{grammarlens}

\subsection*{Your turn}
\begin{itemize}
\raggedright
  \item \emph{Say it:} \emph{cauka}
  \item \emph{Say it:} \emph{cauka}, once more
  \item \emph{Say it:} say \emph{cauka}
  \item \emph{Recall:} say the Kannada for a shape, then the Kannada for a circle, then say \emph{cauka} again
\end{itemize}

\begin{tcolorbox}[breakable,colback=teal!4,colframe=teal!35!black,title={Before you move on}]
What does \kn{ಚೌಕ} mean? (\textquotedblleft{}A square\textquotedblright{}.) Say it once more.
\end{tcolorbox}
\section[\texorpdfstring{\kn{ತ್ರಿಕೋನ} (trikōna)}{ತ್ರಿಕೋನ (trikōna)}]{\kn{ತ್ರಿಕೋನ} (trikōna) — a triangle}

\label{lesson:KA-C193-trikona}



\begin{quote}
Before the new one: say the Kannada for a circle, then the Kannada for a square.
\end{quote}

\subsection*{You'll want to know: \kn{ತ್ರಿಕೋನ}}
\textbf{\kn{ತ್ರಿಕೋನ}} — \emph{trikōna} — \textquotedblleft{}a triangle\textquotedblright{}.

\begin{grammarlens}[title={What you've built}]
One more word for this chapter.
\end{grammarlens}

\subsection*{Your turn}
\begin{itemize}
\raggedright
  \item \emph{Say it:} \emph{trikōna}
  \item \emph{Say it:} \emph{trikōna}, once more
  \item \emph{Say it:} say \emph{trikōna}
  \item \emph{Recall:} say the Kannada for a circle, then the Kannada for a square, then say \emph{trikōna} again
\end{itemize}

\begin{tcolorbox}[breakable,colback=teal!4,colframe=teal!35!black,title={Before you move on}]
What does \kn{ತ್ರಿಕೋನ} mean? (\textquotedblleft{}A triangle\textquotedblright{}.) Say it once more.
\end{tcolorbox}
\section[(dialogue)]{Words from chapters 185-188 — a first pass}

\label{lesson:KA-R193-first-pass-words-from-chapters-185-188}



\begin{quote}
Say the words for a square and a triangle.
\end{quote}

\subsection*{Your turn}
Say these aloud:

\begin{itemize}
\raggedright
  \item a banknote, a coin, a discount, a kilogram, a purse
  \item a message, an email, a pencil, a file, a document
  \item a number, a rainbow, snow, shade, darkness
  \item a nest, a creeper, a bud, a thorn, holy basil
\end{itemize}

\begin{tcolorbox}[breakable,colback=teal!4,colframe=teal!35!black,title={Before you move on}]
A banknote? (\textbf{\kn{ನೋಟು}}, \emph{nōṭu}.) A coin? (\textbf{\kn{ನಾಣ್ಯ}}, \emph{nāṇya}.) A discount? (\textbf{\kn{ರಿಯಾಯಿತಿ}}, \emph{riyāyiti}.) A kilogram? (\textbf{\kn{ಕಿಲೋ}}, \emph{kilō}.) A purse? (\textbf{\kn{ಪರ್ಸ್}}, \emph{pars}.) A message? (\textbf{\kn{ಸಂದೇಶ}}, \emph{sandēśa}.) An email? (\textbf{\kn{ಇಮೇಲ್}}, \emph{imēl}.) A pencil? (\textbf{\kn{ಪೆನ್ಸಿಲ್}}, \emph{pensil}.) A file? (\textbf{\kn{ಕಡತ}}, \emph{kaḍata}.) A document? (\textbf{\kn{ದಾಖಲೆ}}, \emph{dākhale}.) A number? (\textbf{\kn{ಸಂಖ್ಯೆ}}, \emph{saṅkhye}.) A rainbow? (\textbf{\kn{ಕಾಮನಬಿಲ್ಲು}}, \emph{kāmanabillu}.) Snow? (\textbf{\kn{ಹಿಮ}}, \emph{hima}.) Shade? (\textbf{\kn{ನೆರಳು}}, \emph{neraḷu}.) Darkness? (\textbf{\kn{ಕತ್ತಲೆ}}, \emph{kattale}.) A nest? (\textbf{\kn{ಗೂಡು}}, \emph{gūḍu}.) A creeper? (\textbf{\kn{ಬಳ್ಳಿ}}, \emph{baḷḷi}.) A bud? (\textbf{\kn{ಮೊಗ್ಗು}}, \emph{moggu}.) A thorn? (\textbf{\kn{ಮುಳ್ಳು}}, \emph{muḷḷu}.) Holy basil? (\textbf{\kn{ತುಳಸಿ}}, \emph{tuḷasi}.)
\end{tcolorbox}
\section[(dialogue)]{Words from chapters 189-193 — a second pass}

\label{lesson:KA-R193-second-pass-words-from-chapters-189-193}



\begin{quote}
Say the words for a group and a village fair.
\end{quote}

\subsection*{Your turn}
Say these aloud:

\begin{itemize}
\raggedright
  \item a group, a village fair, worship, god, fasting
  \item a sofa, a bench, a stool, a mattress, an axe
  \item a picture, sport, a match, a team, cricket
  \item a grinding stone, a pestle, a mixer-grinder, a sieve, a coconut grater
  \item cardboard, a shape, a circle, a square, a triangle
\end{itemize}

\begin{tcolorbox}[breakable,colback=teal!4,colframe=teal!35!black,title={Before you move on}]
A group? (\textbf{\kn{ಗುಂಪು}}, \emph{gumpu}.) A village fair? (\textbf{\kn{ಜಾತ್ರೆ}}, \emph{jātre}.) Worship? (\textbf{\kn{ಪೂಜೆ}}, \emph{pūje}.) God? (\textbf{\kn{ದೇವರು}}, \emph{dēvaru}.) Fasting? (\textbf{\kn{ಉಪವಾಸ}}, \emph{upavāsa}.) A sofa? (\textbf{\kn{ಸೋಫಾ}}, \emph{sōphā}.) A bench? (\textbf{\kn{ಬೆಂಚು}}, \emph{beñcu}.) A stool? (\textbf{\kn{ಸ್ಟೂಲು}}, \emph{sṭūlu}.) A mattress? (\textbf{\kn{ಹಾಸಿಗೆ}}, \emph{hāsige}.) An axe? (\textbf{\kn{ಕೊಡಲಿ}}, \emph{koḍali}.) A picture? (\textbf{\kn{ಚಿತ್ರ}}, \emph{citra}.) Sport? (\textbf{\kn{ಕ್ರೀಡೆ}}, \emph{krīḍe}.) A match? (\textbf{\kn{ಪಂದ್ಯ}}, \emph{pandya}.) A team? (\textbf{\kn{ತಂಡ}}, \emph{taṇḍa}.) Cricket? (\textbf{\kn{ಕ್ರಿಕೆಟ್}}, \emph{krikeṭ}.) A grinding stone? (\textbf{\kn{ಒರಳು}}, \emph{oraḷu}.) A pestle? (\textbf{\kn{ಒನಕೆ}}, \emph{onake}.) A mixer-grinder? (\textbf{\kn{ಮಿಕ್ಸಿ}}, \emph{miksi}.) A sieve? (\textbf{\kn{ಜರಡಿ}}, \emph{jaraḍi}.) A coconut grater? (\textbf{\kn{ತುರಿಮಣೆ}}, \emph{turimaṇe}.) Cardboard? (\textbf{\kn{ರಟ್ಟು}}, \emph{raṭṭu}.) A shape? (\textbf{\kn{ಆಕಾರ}}, \emph{ākāra}.) A circle? (\textbf{\kn{ವೃತ್ತ}}, \emph{vṛtta}.) A square? (\textbf{\kn{ಚೌಕ}}, \emph{cauka}.) A triangle? (\textbf{\kn{ತ್ರಿಕೋನ}}, \emph{trikōna}.)
\end{tcolorbox}
